\section{CM sources of the theta one-forms}\label{sec:cm-sources}

We now prove Theorem~\ref{thm:cm-extraction}, which identifies the
rational CM sources of the theta classes.  Section~\ref{sec:arithmetic}
constructed an algebraic Hecke action and a cubic relation with
Frobenius.  We calculate the active Hecke characters of the theta
classes, insert them into that relation, and use degree-one Weil
weight to isolate one Frobenius slope on every full character space.
The Hecke projectors preserve the Hodge filtration, so these slopes
determine the Hodge types of all scalar conjugates.  Scalar descent
will then give the required rational Hodge maps.

Retain the fields, moduli space, Albanese product, and Hecke operators
of Section~\ref{sec:arithmetic}.  In particular, the split good prime
\(q\) was chosen after fixing the finite theta inputs and their
common principal level.

\subsection{The active spherical parameter}\label{subsec:active-parameter}

We calculate the active Hecke characters in two steps.  First we
identify the geometric Hecke action on theta inputs as the inverse
spherical action.  We then compute that action on local inputs fixed
by \(\SL_3(\Z_q)\), which form a cyclic module for a dilation
operator.  Rational scalar covariance will
make its eigenvalues algebraic and encode the compact-place signs in
their \(q\)-adic valuations.  A lattice count then identifies the
spherical parameter, with the raw Hecke normalization fixed.

Let \(\mathcal H_w^0\) be the rational subalgebra supported on
positions with multiplier exponent zero, every nonactive selected
block trivial, and active exponent sum zero.  It is the spherical
algebra of \(\SL_3(\Q_q)\).

\begin{lemma}\label{lem:active-hecke-theta}
Every double coset in \(\mathcal H_w^0\) has all its branches
within each component of \(S\).  On the original ball component,
its pullback action on theta classes is the inverse spherical
action on the active local Weil input: a local double coset acts
through its inverse.
\end{lemma}

\begin{proof}
As a quasi-isogeny, such a position leaves the rational polarized
period fixed.  It changes only the active lattice and its opposite
dual, keeping the marking and all other local data.  This is also
the description obtained by clearing the central shift in
Lemma~\ref{lem:hecke-albanese} and undoing its marking change.
A determinant-exponent-zero lattice position has a representative
in \(\SL_3(\Q_q)\).  Adjusting determinant units in Smith bases
shows that its double cosets and branches are exactly those for
\(\SL_3(\Z_q)\backslash\SL_3(\Q_q)/\SL_3(\Z_q)\).

For each component apply strong approximation to the
special-unitary group of its rational Hermitian space.  This
group is simply connected and absolutely almost simple over \(F\),
and its factor at \(\mu\) is noncompact.  Strong approximation
away from all infinite places therefore makes its \(F\)-points
dense in its finite adelic points
\cite[Theorem~4.16]{MilneShimura2017}.
Put \(K_q=\SL_3(\Z_q)\).  Write \(\Lambda^{(j)}\) for the marked
lattice of the component under consideration.  For a branch with
new active lattice \(g^{-1}\Lambda_w^{(j)}\), write
\((\Lambda^{(j)})'\) for the full changed lattice and choose a rational representative
\(\gamma\) whose active component lies in \(K_qg\) and whose other
finite components are integral and marking-preserving.  Then
\(\gamma(\Lambda^{(j)})'=\Lambda^{(j)}\), and \(\gamma\) identifies
the target tuple with the transformed negative period line in the
same component.

On the original component, Lemmas~\ref{lem:theta-tensor-law}
and~\ref{lem:unitary-splitting} give
\[
\gamma^*\Theta\bigl(\omega_f(\gamma)\phi\bigr)=\Theta(\phi).
\]
Thus fixed-input pullback by the branch acts through
\(\omega_f(\gamma)^{-1}\).  Away from \(w\) the chosen level fixes
the input.  At \(w\) our good-prime choice makes it invariant
under \(K_q\); if \(\gamma_w=kg\), its inverse action on that
input is therefore \(\omega_f(g^{-1})\).  Summing the branches
proves the inverse spherical formula.  The sum is raw: no
division by its degree and no additional power of \(q\) occurs.
\end{proof}

We now compute this inverse spherical action on the local inputs.

At \(w\), choose the Schrödinger position space to be the
three-dimensional \(w\)-embedding summand, with lattice \(\Z_q^3\).
By Lemma~\ref{lem:unitary-splitting}, the action of
\(\SL_3(\Q_q)\) is inverse linear substitution.  Every finite input
is a finite sum of tensors separating this place from the other
finite places.  Our choice of \(q\) makes the input invariant under
\(\SL_3(\Z_q)\); averaging a separated decomposition over that
compact group gives one whose local factors are invariant.

The invariant Schwartz functions are radial: \(\SL_3(\Z_q)\) is
transitive on vectors with a fixed minimum coordinate valuation.
Local constancy at zero and compact support therefore make this
space the span of
\[
b_j=\mathbf 1_{q^j\Z_q^3},\qquad j\in\Z.
\]
The functions \(b_j\) are linearly independent.  Define the
normalized dilation
\begin{equation}\label{eq:radial-dilation}
(\mathcal D_q b)(x)=q^{3/2}b(q^{-1}x).
\end{equation}
Then \(\mathcal D_qb_j=q^{3/2}b_{j+1}\), so the radial space is a
cyclic module for \(\C[\mathcal D_q,\mathcal D_q^{-1}]\).

\begin{lemma}\label{lem:scalar-radial-relation}
Fix the input at all finite places other than \(w\).  There are
\(M>0\), an element \(r\in E^\times\), and a nonzero algebraic
number \(\omega\) such that
\[
r\bar r=1,\qquad
\operatorname{val}_{w}(r)=M,\qquad
\operatorname{val}_{\bar w}(r)=-M,
\]
with zero valuation at every other finite prime, and the map from
radial inputs to theta one-form classes factors through
\(\mathcal D_q^M=\omega\).  In the covariance convention of
Lemma~\ref{lem:theta-tensor-law},
\begin{equation}\label{eq:radial-scalar-weight}
\omega^2=\zeta\prod_{\tau\in\Sigma}\tau(r)^{-3d(\tau)}
\end{equation}
for a root of unity \(\zeta\).  Every root of \(X^M-\omega\) has
absolute value one under every complex embedding.
\end{lemma}

\begin{proof}
Principalize a positive power \(w^m=(a)\) and put \(r=a/\bar a\).
It has norm one and the required valuations with \(M=m\).
At every other finite place it is a unit.  Powers of an element of
a compact unit group enter every prescribed open subgroup: its
image in the finite quotient by that subgroup has finite order.
Smoothness of the bare local Weil representation gives an open
stabilizer of the fixed input in the metaplectic group; its image
meets the scalar unit group in an open subgroup.
Replacing \(r\) by a common positive power therefore places it
in the stabilizers needed at the finitely many exceptional places
of the fixed input; replace \(M\) accordingly.

Multiplication by \(r\) on \(V\otimes_E P\) preserves its
\(F\)-symplectic form because \(r\bar r=1\).  We use its
\emph{bare rational symplectic} covariance, with paired
metaplectic lifts, not the value of a chosen unitary splitting
character.  For a transformation \(h\) of a local position space,
the normalized linear action is
\[
(\rho(h)b)(x)=|\det h|_q^{-1/2}b(h^{-1}x),
\]
up to the finite Weil phase of the bare lift.  The linear-subgroup
formula, rational summation invariance, and the double-cover
control are in
\cite[I, no.~13; III, no.~41; IV, nos.~43--45]{Weil1964}.
At \(w\), the position action is the scalar \(r_w\), of valuation
\(M\).  Its unit part fixes a radial function, so the finite
action there is \(\mathcal D_q^M\), up to a root of unity.
At almost every other place the input is the unramified lattice
vector and \(r\) is integral.  At the remaining places the
chosen power fixes the input under a suitable bare lift;
the paired lift differs only by its finite phase.  Their product
is again a root of unity.  An arbitrary lift in the circle
pushout would not justify this conclusion, which is why the bare
rational action is used here.

At infinity, multiplication by \(r\) acts on the \(\tau\)-embedding
summand as the scalar \(\tau(r)\): this is the defining scalar
action after tensoring the \(E\)-space through \(\tau\).
On the \(c\tau\)-summand it acts as
\(\tau(\bar r)=\tau(r)^{-1}\).  It consequently fixes every
chosen positive Lagrangian and the point of \(D\).  In the
identity \(g e_Z=e_{gZ}J\) of
Lemma~\ref{lem:theta-tensor-law}, \(J\) is this action on the
homology filtration.  The squared scalar on the projected first
tensor is the scalar on \(\det\mathcal F_P\) times the square of
the scalar on its projected \(c\)-summand.
At \(\mu\), the filtration has scalar exponents \(1,-1,-1\).
Its determinant contributes \(-1\), and the squared projected
conjugate factor contributes \(-2\).  The total is \(-3\).
At a compact pair, the selected full rank-three summand
contributes \(-3d(\tau)\), with no projected tangent factor.
Thus the squared infinite scalar is
\(\prod_{\tau\in\Sigma}\tau(r)^{-3d(\tau)}\).
This calculation includes the scalar action on the determinant's
constant line, which is trivial only on \(\SU(V)\).

The covariance formula has the finite transformed input on its
left and the tensor action on its right.  Since \(r\) fixes the
period, it gives
\(\Theta(\mathcal D_q^M b)=\omega\Theta(b)\), with the finite
phase absorbed in \(\omega\), and its square is exactly
\eqref{eq:radial-scalar-weight}.  This equality of one-forms
implies the asserted equality of cohomology classes.

The right side of \eqref{eq:radial-scalar-weight} is algebraic,
so \(\omega\) is algebraic.  Every complex embedding of \(E\)
sends \(r\) to a number of modulus one, since \(E\) is CM and
\(r\bar r=1\).  Applying any scalar conjugation to the displayed
identity shows that every complex conjugate of \(\omega\) has
modulus one.  The same is then true of each root of \(X^M-\omega\).
\end{proof}

For a finite separated sum of inputs, take the common stabilizers
of all its other-place factors in this proof.  The same \(r,M,\omega\)
then works for that sum.

We use the half-modulus Satake convention, namely the parameter
attached to normalized parabolic induction.  We do not apply the
additional central Hecke/Tate shift of Gross's Section~8.  In particular the
trivial representation of \(\GL_2\) has parameter
\(\{q^{1/2},q^{-1/2}\}\).  On \(\SL_3\), a parameter is an unordered
triple of nonzero scalars modulo a common scalar.
We say that an active character occurs in a theta class when its
simultaneous generalized eigenspace in the finite-dimensional
Hecke span has a nonzero projection of that class.

\begin{proposition}\label{prop:theta-satake}
Every active spherical character occurring in a theta class of
the multiplicity line \(P\) has half-modulus projective parameter
\begin{equation}\label{eq:active-satake-triple}
\{q^{1/2},q^{-1/2},e\},
\end{equation}
where \(e=e_0^{-1}\) for a root \(e_0^M=\omega\) from
Lemma~\ref{lem:scalar-radial-relation}.  It is algebraic and all
its complex conjugates have modulus one.  For every
\(\sigma\in\Gal(\Qbar/\Q)\),
\begin{equation}\label{eq:active-conjugate-valuations}
\nu_\iota(\sigma e)
 =-\frac32d\bigl((\sigma|_E)^{-1}\bigr)
 =-\frac32v(\sigma|_E).
\end{equation}
A scalar-conjugate active character has a representative of the
form \eqref{eq:active-satake-triple} with third entry
\(e'=\pm\sigma e\).
\end{proposition}

\begin{proof}
The active spherical action commutes with \(\mathcal D_q\).
The radial Laurent module modulo \(\mathcal D_q^M-\omega\)
is a direct sum of one-dimensional spaces
\(\mathcal D_q=e_0\), because \(X^M-\omega\) has distinct
nonzero roots.  It suffices to compute the spherical character
on one of them.  This decomposition concerns the radial
quotient only; it asserts no semisimplicity for the ambient
cohomological Hecke action.

For this computation, extend the substitutions to
\[
\rho(g)b(x)=|\det g|_q^{-1/2}b(g^{-1}x),
\qquad g\in\GL_3(\Q_q),
\]
and put \(K_3=\GL_3(\Z_q)\).  This is a computational extension
of the active \(\SL_3\)-action, not a claim of theta equivariance
for the full moduli similitude group.  Let \(\mathsf T_i\) be
the raw double coset of
\(\diag(q,\ldots,q,1,\ldots,1)\), with \(i\) entries \(q\),
and \(\operatorname{vol}(K_3)=1\).  Its right cosets correspond
to sublattices between \(q\Z_q^3\) and \(\Z_q^3\) of index
\(q^i\).  For \(i=1,2\) there are \(q^2+q+1\) such lattices.
The numbers containing a fixed primitive vector are respectively
\(q+1\) and \(1\).  Since \(\rho(g)b_0\) is
\(q^{i/2}\) times the characteristic function of its sublattice,
we obtain
\begin{equation}\label{eq:radial-minuscule-actions}
\begin{aligned}
\mathsf T_1b_0&=q^{1/2}\bigl((q+1)b_0+q^2b_1\bigr),\\
\mathsf T_2b_0&=q\bigl(b_0+(q^2+q)b_1\bigr),\\
\mathsf T_3b_0&=\mathcal D_qb_0.
\end{aligned}
\end{equation}
For example, vectors in \(q\Z_q^3\) belong to every sublattice,
while a primitive vector is tested by its nonzero class in
\(\Fq^3\); the displayed counts are the numbers of the
corresponding planes or lines containing that class.

On the quotient \(\mathcal D_q=e_0\) we have
\(b_1=q^{-3/2}e_0b_0\).  Substitution into
\eqref{eq:radial-minuscule-actions} gives the values
\[
q\,e_1(e_0,q^{1/2},q^{-1/2}),\qquad
q\,e_2(e_0,q^{1/2},q^{-1/2}),\qquad
e_3(e_0,q^{1/2},q^{-1/2}),
\]
where \(e_i\) is the \(i\)-th elementary symmetric polynomial.
In the half-modulus Satake isomorphism the raw minuscule cosets
have transforms \(q e_1,q e_2,e_3\); these generate the spherical
algebra, with \(e_3\) invertible.  See
\cite[Equation~(3.14) and \S6]{Gross1998}, where the
\(\GL_n\) formula is attributed to Tamagawa.  The character just
computed therefore has parameter
\(\{e_0,q^{1/2},q^{-1/2}\}\).

The geometric action on theta inputs is inverse by
Lemma~\ref{lem:active-hecke-theta}.  Inversion sends the
half-modulus parameter to its reciprocal, without a degree
factor.  One can check this on the same generators:
\(\mathsf T_1^\vee=\mathsf T_3^{-1}\mathsf T_2\),
\(\mathsf T_2^\vee=\mathsf T_3^{-1}\mathsf T_1\), and
\(\mathsf T_3^\vee=\mathsf T_3^{-1}\), where \(f^\vee(g)=f(g^{-1})\).
Their transforms are the corresponding elementary symmetric
functions of the reciprocal triple.  After interchanging the
first two entries, the active projective parameter is therefore
\(\{q^{1/2},q^{-1/2},e_0^{-1}\}\).

The determinant-zero \(\GL_3\) and \(\SL_3\) double cosets agree
as in Lemma~\ref{lem:active-hecke-theta}.  Their Satake unipotent
integrals and modulus factors agree as well.  The half-sum of
positive roots is \((1,0,-1)\), integral on the relevant
cocharacter lattices; the Satake comparison is consequently
defined over rational coefficients
\cite[Proposition~3.6]{Gross1998}.  Scalar conjugation of an
active numerical character therefore conjugates its projective
parameter.  It may change the two entries \(q^{1/2},q^{-1/2}\)
by their common sign.  Removing that sign projectively leaves
third entry \(\pm\sigma e\).

It remains to calculate the valuation.  Put \(g=\sigma|_E\).
By the definition of the prime labels,
\[
\nu_\iota(\sigma\tau(r))
 =\operatorname{val}_{\mathfrak p_{g\tau}}(r).
\]
This is nonzero only when \(g\tau=1\) or \(c\).  Among the
representatives \(\Sigma\), the contributing one is \(g^{-1}\)
with valuation \(-M\), or \(cg^{-1}\) with valuation \(M\).
Here \(c\) is central in \(G_E\), and oddness of \(d\) gives
in either case
\[
\sum_{\tau\in\Sigma}d(\tau)\nu_\iota(\sigma\tau(r))
   =-M d(g^{-1}).
\]
Taking valuations in \eqref{eq:radial-scalar-weight} and in
\(e_0^M=\omega\) yields
\(\nu_\iota(\sigma e_0)=\tfrac32d(g^{-1})\).
Since \(e=e_0^{-1}\) and \(d(g^{-1})=v(g)\), this proves
\eqref{eq:active-conjugate-valuations}.  The unit-modulus
assertion follows from Lemma~\ref{lem:scalar-radial-relation}.
\end{proof}

The active character is now known, including the valuation of every
scalar conjugate.  The full Hecke character may still multiply its
three Satake entries by a common scalar.  The following finite-order
relations will control that scalar when we apply the Frobenius
polynomial of Proposition~\ref{prop:albanese-congruence}.

\begin{lemma}\label{lem:central-finite-order}
The central operators \(T_3\) and \(C_0\) have finite order on
the full tuple space \(S\) and on \(A_S\).
\end{lemma}

\begin{proof}
Each central position has a unique lattice choice.  For \(T_3\),
use the norm-one element \(r\) constructed in
Lemma~\ref{lem:scalar-radial-relation}; only its valuations and
norm are needed here.  The scalar \(r\) takes the \(T_3^M\)-target
lattice back to the source lattice on the active and opposite
blocks.  It preserves every polarized period
because it is an \(E\)-scalar of norm one.  Its denominators are
supported at \(q\), so a further power is one on the \(N\)-marking.
That power of \(T_3\) is the identity on every tuple.

For \(C_0\), let \(I\) be the product of the primes where its
full exponent is one.  There is exactly one from each conjugate
pair, including \(\bar w\), and \(I\bar I=(q)\).  Choose
\(I^m=(a)\).  Then
\[
a\bar a=q^m u
\]
for a totally positive unit \(u\in\mathcal O_F^\times\).
The scalar \(b=a^2/u\) has exactly the valuations of \(I^{2m}\)
and satisfies
\[
b\bar b=q^{2m}.
\]
It identifies the corresponding powered target lattice with
the source, with exactly the required polarization multiplier.
A further power fixes the marking.  This argument works on
every rational Hermitian component because an \(E\)-scalar
preserves each of its embedding filtrations.  Thus a power of
\(C_0\) is also the identity on the full tuple space.  The
induced Albanese homomorphisms have finite order as well.
\end{proof}
